%%
%% ACM TODAES Journal Paper
%% Special Issue: Open, Reliable, and Generalizable Datasets for AI in EDA
%% Submission deadline: August 31, 2026
%%
%% This paper substantially extends our workshop paper:
%%   CTS-Bench: Benchmarking Graph Coarsening Trade-offs for GNNs in Clock
%%   Tree Synthesis (ML Bench'26 @ ASPLOS, 2026)
%% New contributions: 19-design corpus, HPC pipeline, ~10K dataset,
%%   SwiftCTS surrogate as proof of dataset utility.
%%

\documentclass[acmsmall,review,anonymous]{acmart}

\usepackage{booktabs}
\usepackage{multirow}
\usepackage{amsmath,amssymb}
\usepackage{graphicx}
\usepackage{enumitem}
\usepackage{xcolor}
\usepackage{colortbl}
\usepackage[table]{xcolor}
\usepackage{soul}
\usepackage{tikz}
\usetikzlibrary{shapes,arrows,positioning,fit,backgrounds,calc,shadows,
                shapes.geometric,shapes.symbols,shapes.misc,patterns}

\definecolor{bestrow}{RGB}{213,240,213}

%% Journal metadata
\acmJournal{TODAES}
\acmVolume{0}
\acmNumber{0}
\acmArticle{0}
\acmYear{2026}
\acmMonth{8}
\copyrightyear{2026}
\acmDOI{}

\begin{document}

\title{CTS-Bench: An Open, Large-Scale Dataset and Automated Pipeline for
       AI-Driven Clock Tree Synthesis}

\author{Barsat Khadka}
\affiliation{%
  \institution{The University of Southern Mississippi}
  \city{Hattiesburg}
  \state{Mississippi}
  \country{USA}}
\email{Barsat.Khadka@usm.edu}

\author{Kawsher Ahmed Roxy}
\affiliation{%
  \institution{Intel Corporation}
  \city{Hillsboro}
  \state{Oregon}
  \country{USA}}
\email{kawsher.roxy@intel.com}

\author{Md Rubel Ahmed}
\affiliation{%
  \institution{Louisiana Tech University}
  \city{Ruston}
  \state{Louisiana}
  \country{USA}}
\email{mahamed@latech.edu}

\begin{abstract}
Machine learning for Clock Tree Synthesis (CTS) remains fundamentally
constrained by the lack of open, large-scale, and reproducible datasets.
Existing efforts use at most five proprietary or limited-license designs
and rarely expose the complete placement-to-CTS pipeline, preventing
community reuse. We present \textbf{CTS-Bench}, a dataset and reproducible
generation framework that systematically bridges this gap.

CTS-Bench provides \textbf{10,000+ CTS evaluations} derived from
\textbf{19 open-source RTL designs} spanning two diverse benchmark suites
(14 IWLS and 5 PlaceDreamer) across the SkyWater 130\,nm PDK.
Each design undergoes randomized placement under seven knobs followed by
ten distinct CTS configurations, capturing how structural placement
decisions propagate into clock power, timing skew, wirelength, and routing
complexity. An automated, containerized HPC pipeline built on
OpenLane~\cite{OpenRoad}, Singularity/Apptainer, and SLURM ensures
full end-to-end reproducibility and allows the benchmark to be extended to
new designs or technology nodes with a single command.

Beyond the dataset itself, we demonstrate its utility along two axes.
First, we provide multi-scale graph representations (raw gate-level and
physics-aware clustered graphs) enabling GNN benchmarking studies;
experiments with GCN, GraphSAGE, and GATv2 reveal that graph coarsening
reduces memory by up to $17.2\times$ but breaks local clock skew
prediction, motivating task-aware clustering research.
Second, we introduce \textbf{SwiftCTS}, a physics-informed surrogate
trained entirely on CTS-Bench data, which predicts CTS outcomes in
sub-millisecond latency---enabling evaluation of 100,000 configurations in
under 10 seconds---and reduces power prediction error from 24.5\% to 3.3\%
on completely unseen macro architectures with only one physical reference
run. Together, CTS-Bench and SwiftCTS establish a foundation for
data-driven CTS research: an open, large-scale dataset and an end-to-end
demonstration of what becomes possible when such data exists.

All data, code, and the HPC pipeline are publicly available at:
\href{https://github.com/BarsatKhadka/SwiftCTS}{https://github.com/BarsatKhadka/SwiftCTS}
\end{abstract}

\begin{CCSXML}
<ccs2012>
  <concept>
    <concept_id>10010583.10010682.10010697.10010698</concept_id>
    <concept_desc>Hardware~Clock-network synthesis</concept_desc>
    <concept_significance>500</concept_significance>
  </concept>
  <concept>
    <concept_id>10010583.10010682.10010697.10010701</concept_id>
    <concept_desc>Hardware~Placement</concept_desc>
    <concept_significance>500</concept_significance>
  </concept>
  <concept>
    <concept_id>10010147.10010257.10010293.10010294</concept_id>
    <concept_desc>Computing methodologies~Neural networks</concept_desc>
    <concept_significance>300</concept_significance>
  </concept>
  <concept>
    <concept_id>10010583.10010682.10010712.10010713</concept_id>
    <concept_desc>Hardware~Best practices for EDA</concept_desc>
    <concept_significance>300</concept_significance>
  </concept>
</ccs2012>
\end{CCSXML}

\ccsdesc[500]{Hardware~Clock-network synthesis}
\ccsdesc[500]{Hardware~Placement}
\ccsdesc[300]{Computing methodologies~Neural networks}
\ccsdesc[300]{Hardware~Best practices for EDA}

\keywords{Electronic Design Automation, Clock Tree Synthesis, Open Dataset,
  Machine Learning, Graph Neural Networks, Surrogate Models, Physical Design,
  HPC Pipeline, OpenROAD, Sky130}

\maketitle

%%────────────────────────────────────────────────────────────────────────────
\section{Introduction}
\label{sec:intro}
%%────────────────────────────────────────────────────────────────────────────

Clock Tree Synthesis is among the most computationally intensive and
quality-sensitive stages of the physical design flow. The clock network
directly governs chip power (15--40\% of total dynamic dissipation in
modern designs), timing closure, and routing congestion. Engineers
navigate a high-dimensional configuration space---cluster diameter, maximum
wire length, buffer distances, and sink cluster sizes---by repeatedly
invoking the EDA tool on each placement, an exhaustive process that can
consume days per tapeout. The EDA community has turned increasingly to
machine learning (ML) to automate this exploration, proposing analytical
models~\cite{kahng,kahng2,kahng3}, convolutional
surrogates~\cite{nagaria2020}, generative adversarial
frameworks~\cite{gan-cts}, and graph neural
networks~\cite{wu2025graph,Hu2023SyncTREE}, as well as reinforcement
learning and Bayesian optimization~\cite{agnesina,thomas2025,geng2022,zheng2023}.

Despite this progress, every proposed method faces the same fundamental
bottleneck: the absence of a public, large-scale, and reproducible
dataset tailored to the placement-CTS interface.

\paragraph{The data problem.}
Existing studies train and evaluate on two to six proprietary or
restricted-license designs \cite{gan-cts,nagaria2020}, generate at most a
few thousand evaluations, and provide neither the data nor the generation
pipeline. As the NSF Workshop on AI for EDA explicitly
concluded~\cite{nsfReport}, poor generalization in ML-EDA is largely a
data problem: models trained on limited design families collapse when faced
with unseen macro architectures \cite{zhengDAC,ghoseICCAD}, and without
public datasets, every new research group must rebuild their own
infrastructure from scratch.
The TODAES Special Issue on Open, Reliable, and Generalizable Datasets for
AI in EDA was created precisely because ``design and manufacturing data are
fragmented across abstraction layers, tools, and technology nodes, and
existing datasets often suffer from limited scale, quality, accessibility,
and diversity.''
CTS-Bench directly addresses this gap.

\paragraph{Contributions.}
This paper makes four primary contributions:
\begin{enumerate}[leftmargin=*,topsep=2pt,itemsep=2pt]
  \item \textbf{Scale and diversity.} We release a dataset of \textbf{10,000+
    CTS evaluations} across \textbf{19 open-source RTL designs} (14 IWLS + 5
    PlaceDreamer), spanning a $100\times$ range in design complexity, all
    synthesized and placed under the SkyWater 130\,nm PDK with OpenROAD.

  \item \textbf{Reproducible HPC pipeline.} We provide a fully automated,
    containerized generation pipeline (OpenLane inside Singularity/Apptainer,
    SLURM array jobs) that can be re-run or extended to new designs or PDKs
    with a single command. Every placement configuration is logged and every
    CTS run archived.

  \item \textbf{Multi-scale graph representations.} Each physical placement is
    accompanied by both a raw gate-level graph and a physics-aware clustered
    graph. GNN benchmarking experiments on GCN, GraphSAGE, and GATv2 quantify
    the accuracy--efficiency trade-off and reveal that generic graph coarsening
    breaks local clock skew prediction.

  \item \textbf{Proof of dataset utility: SwiftCTS.} Trained entirely on
    CTS-Bench data, SwiftCTS is a physics-informed surrogate (gradient-boosted
    ensembles + K-shot calibration) that predicts CTS outcomes at
    sub-millisecond latency, enables comprehensive Pareto optimization over
    100,000 configurations in under 10 seconds, and generalizes to unseen
    macro architectures with under 3.3\% power error after a single physical
    reference run.
\end{enumerate}

\noindent
This work substantially extends our preliminary workshop
paper~\cite{khadka2026} (5 designs, 4,860 evaluations, no HPC pipeline,
no surrogate). New material comprises the 19-design corpus, the
production-grade HPC pipeline, the SwiftCTS surrogate and its complete
experimental evaluation, and an expanded dataset characterization with
cross-design statistical analysis.

The rest of this paper is organized as follows.
Section~\ref{sec:related} reviews related work.
Section~\ref{sec:dataset} describes the CTS-Bench dataset, design corpus,
and generation pipeline.
Section~\ref{sec:graphs} covers multi-scale graph representations.
Section~\ref{sec:analysis} characterizes the dataset.
Section~\ref{sec:swiftcts} presents the SwiftCTS surrogate.
Section~\ref{sec:gnn} reports GNN benchmarking experiments.
Section~\ref{sec:surrogate_results} presents SwiftCTS results.
Section~\ref{sec:limitations} discusses limitations and future directions.
Section~\ref{sec:conclusion} concludes.

%%────────────────────────────────────────────────────────────────────────────
\section{Background and Related Work}
\label{sec:related}
%%────────────────────────────────────────────────────────────────────────────

\subsection{Machine Learning for CTS}

\textit{Analytical and metamodeling approaches.}
Kahng et al.~\cite{kahng,kahng2,kahng3} pioneered rule-based data mining and
high-dimensional metamodeling for CTS outcome prediction, using features such
as gate count, whitespace, and floorplan aspect ratio. While efficient, these
approaches rely on hand-engineered aggregate features and degrade significantly
on unseen designs, with reported correlations dropping from near-perfect to
0.48 on new configurations.

\textit{Deep learning architectures.}
Nagaria and Deb~\cite{nagaria2020} employ CNNs to estimate CTS parameters
from pre-CTS placement images. Lu et al.~\cite{gan-cts} introduce GAN-CTS, a
generative adversarial framework that exports high-resolution spatial maps and
processes them through a 50-layer convolution, achieving impressive per-design
accuracy but with a 45\% failure rate on the NOVA benchmark and inference
latency far too high for exhaustive search. Recent transfer learning
approaches~\cite{zhang2024,zhao2025} address cross-technology generalization
but remain data-hungry.

\textit{Reinforcement learning and Bayesian optimization.}
RL-based frameworks~\cite{agnesina,thomas2025} use GNNs and Transformers to
tune physical design parameters but require the EDA tool inside the
optimization loop, prohibiting exhaustive search.
Bayesian optimization~\cite{geng2022,zheng2023} samples more efficiently but
still invokes the tool at each evaluation.

The common thread across all existing approaches: they were developed without
a shared dataset. CTS-Bench makes training, evaluation, and cross-method
comparison on a common corpus possible for the first time.

\subsection{EDA Datasets and Benchmarks}

\textit{Traditional benchmarks.}
ISPD98~\cite{10.1145/274535.274546}, ISPD2005~\cite{10.1145/1055137.1055182},
and ICCAD2015~\cite{7372671} serve algorithmic optimization but provide only
one reference solution per instance, offer no ML-ready labels, and predate
modern deep learning requirements.

\textit{Graph-centric datasets.}
CircuitNet~\cite{jiang2024circuitnet} and
EDA-schema~\cite{10.1145/3649476.3658718} provide broader physical design
datasets. CircuitNet 2.0 covers congestion and timing across multiple
industrial designs; EDA-schema introduces a standardized graph format with
960 designs. Neither isolates the placement-CTS interface or provides the
multi-variant CTS evaluations per placement needed to study configuration
sensitivity.

\textit{General graph benchmarks.}
OGB~\cite{hu2021opengraphbenchmarkdatasets} and
LRGB~\cite{dwivedi2023longrangegraphbenchmark} standardize evaluation on
social, citation, and molecular graphs. Their structural properties
(low-fanout, undirected, small diameter) are fundamentally different from
logic netlists, making them poor proxies for EDA workloads.

\noindent
CTS-Bench fills the gap left by all prior work: a large-scale,
CTS-specific dataset paired with multi-scale graph representations, a
reproducible generation pipeline, and a deployed surrogate demonstrating
dataset utility.

%%────────────────────────────────────────────────────────────────────────────
\section{The CTS-Bench Dataset}
\label{sec:dataset}
%%────────────────────────────────────────────────────────────────────────────

\subsection{Design Corpus}

CTS-Bench comprises 19 open-source RTL designs sourced from two established
benchmark suites. Table~\ref{tab:designs} lists all designs along with their
target clock period, clock port, and top module. The corpus spans four orders
of magnitude in logic complexity, from the 1.8K-cell ZipDiv divider core to
the 100K+ cell JPEG encoder.

\begin{table}[h]
\centering
\small
\caption{CTS-Bench Design Corpus (19 designs).}
\label{tab:designs}
\begin{tabular}{llll}
\toprule
\textbf{Design} & \textbf{Suite} & \textbf{Clk (ns)} & \textbf{Top Module} \\
\midrule
usb\_phy    & IWLS & 10.0 & usb\_phy \\
mem\_ctrl   & IWLS & 10.0 & mc\_top \\
jpeg        & IWLS & 10.0 & jpeg\_top \\
wb\_dma     & IWLS & 10.0 & wb\_dma\_top \\
ac97\_ctrl  & IWLS & 10.0 & ac97\_top \\
pci         & IWLS & 10.0 & pci\_bridge32 \\
i2c         & IWLS & 10.0 & i2c\_master\_top \\
spi         & IWLS & 10.0 & simple\_spi \\
tv80        & IWLS & 10.0 & tv80s \\
aes         & IWLS &  7.0 & aes \\
picorv32    & IWLS &  5.0 & picorv32 \\
sha256      & IWLS &  9.0 & sha256 \\
ethmac      & IWLS &  9.0 & eth\_top \\
zipdiv      & IWLS &  5.0 & zipdiv \\
\midrule
salsa20     & PlaceDreamer & 10.0 & salsa20 \\
xtea        & PlaceDreamer & 10.0 & xtea \\
y\_huff     & PlaceDreamer & 10.0 & y\_huff \\
PPU         & PlaceDreamer & 10.0 & PPU \\
usb         & PlaceDreamer & 20.0 & usb \\
\bottomrule
\end{tabular}
\end{table}

\noindent
\textbf{IWLS designs} are canonical open-source IP cores covering diverse
functional domains: cryptography (AES, SHA-256), microprocessors (PicoRV32,
TV80), networking (EthMAC, USB PHY, PCI bridge), arithmetic (ZipDiv),
and bus interfaces (WB-DMA, I2C, SPI, AC97, Memory Controller, JPEG).

\noindent
\textbf{PlaceDreamer designs} introduce additional structural diversity:
XTEA and Salsa20 are lightweight stream ciphers, Y-Huff is a Huffman codec,
PPU is a pixel processing unit, and USB is a full-speed USB controller.
These designs were sourced from the PlaceDreamer benchmark to extend coverage
beyond the canonical IWLS set.

\subsection{Automated Generation Pipeline}

The generation pipeline proceeds in three stages for each design, as
illustrated in Figure~\ref{fig:pipeline}.

\paragraph{Stage 1: Randomized Placement.}
Starting from RTL, we execute synthesis, floorplanning, power network
insertion, and detailed placement using OpenLane inside an
Apptainer/Singularity container encapsulating OpenLane~2.3.10 with the
SkyWater 130\,nm PDK.
We stop immediately before CTS.
To create a diverse placement corpus, we randomize seven parameters
(Table~\ref{tab:knobs}) that alter spatial cell distribution,
synthesis strategy, and routing pressure. Each unique combination of
randomized parameters produces one \emph{placement instance}.

\paragraph{Stage 2: Activity Extraction (SAIF).}
For each placement instance, we perform gate-level simulation using the
Icarus Verilog toolchain (invoked inside the same Singularity container),
supplying design-specific testbenches. The resulting VCD waveform is
converted to SAIF format via a Python-native converter, capturing
per-net switching activity for each unique placement. These SAIF files are
essential inputs to the switching-activity features used by SwiftCTS.

\paragraph{Stage 3: Multi-Variant CTS.}
For each placement, we execute 10 distinct CTS runs under TritonCTS by
uniformly sampling four CTS tuning knobs (Table~\ref{tab:knobs}). Static
timing analysis (STA) follows each CTS run to extract 15 ground-truth QoR
metrics (Table~\ref{tab:metrics}).

\begin{figure}[h]
\centering
% TODO: replace with actual pipeline figure (use TikZ or PDF)
\includegraphics[width=\columnwidth]{placeholder_pipeline.pdf}
\caption{CTS-Bench generation pipeline. Stage~1 randomizes placement;
  Stage~2 extracts gate-level switching activity; Stage~3 runs 10 CTS
  variants per placement and extracts 15 QoR metrics.}
\label{fig:pipeline}
\end{figure}

\subsection{HPC Infrastructure}

Generating 10,000+ CTS evaluations is computationally infeasible on a
single workstation.
The CTS-Bench pipeline is designed for deployment on SLURM-managed HPC
clusters.
The entire tool stack---OpenLane 2.3.10, OpenROAD, TritonCTS, and Icarus
Verilog---is packaged into a single Singularity Image File (SIF), eliminating
host dependency requirements beyond Python 3 and Singularity/Apptainer.
SLURM array jobs (200 tasks, partition \texttt{cpuq}) fan out placement and
CTS evaluations in parallel; each task writes a CSV shard that is merged
post-execution. Setup requires only:
\begin{enumerate}[leftmargin=*,topsep=2pt,itemsep=0pt]
  \item \texttt{bash hpc/setup.sh} -- pulls the SIF and enables the PDK
    (one-time, shared filesystem);
  \item \texttt{sbatch hpc/slurm/run\_array.sbatch} -- submits all jobs.
\end{enumerate}
Full reproducibility is achieved via version-pinned tooling
(OpenLane 2.3.10, PDK hash \texttt{0fe599b2}), and the pipeline
auto-detects the Singularity/Apptainer binary and HPC home path.

\subsection{Randomization Knobs}

\begin{table}[h]
\centering
\small
\caption{Randomization Knobs for Dataset Diversity.}
\label{tab:knobs}
\begin{tabular}{llll}
\toprule
\textbf{Stage} & \textbf{Parameter} & \textbf{Range / Set} & \textbf{Type} \\
\midrule
Synthesis   & Strategy      & \{AREA 0-2, DELAY 0-4\}   & Categorical \\
Floorplan   & Aspect Ratio  & \{0.7, 1.0, 1.4, 2.0\}   & Discrete \\
Floorplan   & IO Mode       & \{Random, Standard\}      & Binary \\
Placement   & Core Util.    & 20\%--70\%                & Continuous \\
Placement   & Target Density & Util + [0.0, 0.20]      & Continuous \\
Placement   & Time Driven   & \{On, Off\}               & Binary \\
Placement   & Routability   & \{On, Off\}               & Binary \\
\midrule
CTS & Sink Max Diameter   & 35--70 $\mu$m             & Integer \\
CTS & Max Wire Length     & 130--280 $\mu$m           & Integer \\
CTS & Cluster Size        & 12--30 sinks              & Integer \\
CTS & Buffer Distance     & 70--150 $\mu$m            & Integer \\
\bottomrule
\end{tabular}
\end{table}

\subsection{Ground-Truth QoR Labels}

\begin{table}[h]
\centering
\small
\caption{Ground-Truth QoR Metrics Captured per CTS Run.}
\label{tab:metrics}
\begin{tabular}{lll}
\toprule
\textbf{Category} & \textbf{Metric} & \textbf{Unit} \\
\midrule
Timing   & Skew (Setup/Hold), Slack, TNS & ns \\
Power    & Total Dynamic \& Static Power & W \\
Physical & Wirelength, Cell Utilization  & $\mu$m, \% \\
Resources & Clock Buffers, Inverters     & integer \\
Congestion & Routing / Repair Buffer Count & integer \\
\bottomrule
\end{tabular}
\end{table}

%%────────────────────────────────────────────────────────────────────────────
\section{Multi-Scale Graph Representations}
\label{sec:graphs}
%%────────────────────────────────────────────────────────────────────────────

Each physical placement in CTS-Bench is accompanied by two graph
representations designed to support a range of GNN architectures and hardware
targets.

\subsection{Raw Gate-Level Graph}

Nodes represent standard cells (flip-flops and combinational logic); edges
represent logical connectivity. We prioritize flip-flops and their one-hop
fan-out logic, as data-path cells immediately surrounding registers exert the
strongest influence on clock tree
behavior~\cite{10.1145/3670474.3685949}.

\paragraph{Node features.}
Three features are assigned per node: (1)~normalized geometric coordinates
$(x, y)$; (2)~a binary cell-type indicator (\texttt{is\_ff}); and
(3)~log-transformed switching activity from SAIF.

\paragraph{Edge weights.}
Edges are weighted by Manhattan distance:
$d(c_i, c_j) = |x_i - x_j| + |y_i - y_j|$.

Raw graphs range from roughly 2,000 nodes (ZipDiv) to over 200,000 nodes
(JPEG encoder), stressing GPU memory even on high-end accelerators.

\subsection{Physics-Aware Clustered Graph}

To enable efficient GNN training on memory-constrained hardware, we develop
a three-step, physics-aware clustering algorithm that achieves an average
$13.3\times$ node compression while preserving structural information
critical for CTS target metrics.

\paragraph{Step I: Atomic cluster formation.}
A BFS starting from each flip-flop claims its immediate combinational fan-out
cone until another flip-flop or an already-claimed gate is encountered.
Flip-flops are randomly shuffled before BFS to eliminate ordering bias.

\paragraph{Step II: High-spread filtering.}
Clusters whose member-cell spatial spread (standard deviation) exceeds 0.05
normalized die units are flagged as high-spread and excluded from further
merging, preserving their unique geometric signatures.

\paragraph{Step III: Gravity-vector-aligned merging.}
Compact clusters are merged into macro-nodes subject to three constraints:
(a)~shared control net (reset or enable); (b)~Manhattan distance between
centroids $<0.05$; and (c)~cosine similarity of gravity vectors $>0.9$,
where the \emph{gravity vector} of a flip-flop is the displacement from its
position to the centroid of its one-hop neighbors, encoding the directional
pull of surrounding logic.

\paragraph{Clustered node features.}
Each macro-node carries 10 aggregated features: geometric centroid $(x,y)$,
spatial spread $(\sigma_x, \sigma_y)$, log-transformed cluster size,
composition counts ($N_\text{FF}$, $N_\text{Logic}$), and aggregated switching
activity (log max, sum, and non-zero toggle counts). Edge weights are the
Manhattan distance between macro-cluster centroids.

Both graph types are serialized as PyTorch Geometric tensors and referenced
in the master CSV manifest.

%%────────────────────────────────────────────────────────────────────────────
\section{Dataset Characterization}
\label{sec:analysis}
%%────────────────────────────────────────────────────────────────────────────

\subsection{Scale and Coverage}

\begin{table}[h]
\centering
\small
\caption{CTS-Bench Dataset Statistics.}
\label{tab:stats}
\begin{tabular}{lrr}
\toprule
\textbf{Statistic} & \textbf{Workshop (2026)} & \textbf{This Work} \\
\midrule
Designs          &  5   & 19 \\
Design suites    &  1   & 2 (IWLS + PlaceDreamer) \\
Unique placements &  486 & $\sim$1{,}000+ \\
CTS evaluations  & 4{,}860 & 10{,}000+ \\
QoR metrics/run  & 15   & 15 \\
Graph types      & 2    & 2 (raw + clustered) \\
PDK              & Sky130 & Sky130 \\
Tool stack       & OpenLane 2.3.10 & OpenLane 2.3.10 \\
Pipeline         & Local & HPC (SLURM array, 200 tasks) \\
\bottomrule
\end{tabular}
\end{table}

\subsection{Pareto Gap Analysis}

To characterize the placement-CTS interaction across heterogeneous
architectures, we define a normalized \emph{Pareto Gap} framework.
Raw metrics (skew in ps, power in W, wirelength in $\mu$m) are not
directly comparable across designs of different scales. For every CTS run $i$
of design $D$ we define the gap vector:
\begin{equation}
  \vec{G}_i = \left[
    \frac{\text{Skew}_i}{\min(\text{Skew})_D},\;
    \frac{\text{Power}_i}{\min(\text{Power})_D},\;
    \frac{\text{WL}_i}{\min(\text{WL})_D}
  \right]
\end{equation}
and the total Pareto distance from the ideal anchor $[1,1,1]$:
\begin{equation}
  D_\text{Pareto} = \sqrt{(G_\text{Skew}-1)^2 + (G_\text{Power}-1)^2 +
                          (G_\text{WL}-1)^2}.
\end{equation}
A high $G_\text{Power}$ coupled with low $G_\text{Skew}$ reveals placements
that require aggressive buffer insertion for timing closure, making
$D_\text{Pareto}$ a useful single scalar for placement quality evaluation.
Figure~\ref{fig:pareto_dist} shows the distribution of $D_\text{Pareto}$
across all 19 designs. % TODO: generate figure after HPC run completes

\begin{figure}[h]
\centering
% TODO: replace with actual violin / box plot from dataset
\includegraphics[width=\columnwidth]{placeholder_pareto.pdf}
\caption{Pareto distance distribution across 19 designs. Lower
  $D_\text{Pareto}$ indicates placements that admit clock trees closer to
  simultaneously optimal in all three objectives.}
\label{fig:pareto_dist}
\end{figure}

\subsection{Cross-Design Metric Variability}

Clock tree power varies by $\sim$90$\times$ across the 19-design corpus
(from ZipDiv at sub-mW to JPEG encoder at $>$200\,mW), while within any
single placement it varies only 10--30\% across CTS knob settings. This
order-of-magnitude between-design vs.\ within-design variance is the core
challenge that motivates the K-shot calibration mechanism in SwiftCTS
(Section~\ref{sec:kshot}).

%%────────────────────────────────────────────────────────────────────────────
\section{SwiftCTS: Proof of Dataset Utility}
\label{sec:swiftcts}
%%────────────────────────────────────────────────────────────────────────────

To demonstrate that CTS-Bench enables practical ML research, we present
SwiftCTS: a physics-informed surrogate framework trained on CTS-Bench data
that predicts clock tree metrics at sub-millisecond latency and drives
comprehensive design space exploration. The three prediction targets are
clock power, wirelength, and timing skew.

\subsection{Feature Engineering}
\label{sec:features}

SwiftCTS extracts a 110-dimensional physical feature space from the
placement DEF, SAIF switching activity file, and pre-CTS timing report.
Features are constructed to be \emph{scale-invariant}: logarithmic
compression handles the dynamic range of circuit sizes, dimensionless
ratios enable learning of universal density patterns, and explicit
interaction terms between layout topology and CTS knobs ground the model
in circuit physics. The full feature taxonomy (Table~\ref{tab:features})
is organized into six groups for the wirelength head and dedicated subsets
for power and skew.

\begin{table}[h]
\centering
\fontsize{7pt}{8pt}\selectfont
\setlength{\tabcolsep}{3pt}
\caption{110-Dimensional Physical Feature Space \& Nomenclature.}
\label{tab:features}
\resizebox{\columnwidth}{!}{%
\renewcommand{\arraystretch}{1.1}
\begin{tabular}{|c|c|c|}
\hline
\textbf{Head} & \multicolumn{2}{c|}{\textbf{Mathematical Formulation}} \\ \hline
\multirow{10}{*}{\textbf{\shortstack{Power\\(20)}}}
 & $\log(1+n_{act}\!\times\!\alpha_{rel}\!\times\!f)$ & $\log(1+n_{ff})$ \\ \cline{2-3}
 & $\alpha_{mean}/\alpha_{max}$            & $\log(1+f_{mux}\!\times\!n_{act})$ \\ \cline{2-3}
 & $\log(1+n_{act}\!\times\!\mu_{ds})$    & $n_{mux}/n_{act}$ \\ \cline{2-3}
 & $\mathrm{frac}(\alpha>2\mu_\alpha)$    & $n_{and\_or}/n_{act}$ \\ \cline{2-3}
 & $\sigma(\text{slack})$                 & $\mathrm{frac}(\text{slack}<0.5\,\text{ns})$ \\ \cline{2-3}
 & $\mu(P_{sig=1})$                       & $\text{slack}_{p10}$ \\ \cline{2-3}
 & $\text{slack}_{p50}$                   & $\mu_\text{slack}\!\times\!f_{xor}$ \\ \cline{2-3}
 & $1/t_{clk}$                            & $\text{Synth}_\text{AreaWeight}$ \\ \cline{2-3}
 & $\log(1+f_{xor}\!\times\!n_{act})$     & $\log(1+n_{nets})$ \\ \cline{2-3}
 & $n_{nand\_nor}/n_{act}$                & $\log(1+cd\!\times\!n_{ff}/\text{Area})$ \\ \hline
\multirow{6}{*}{\textbf{\shortstack{Skew\\(15)}}}
 & $\log(1+\{cd,cs,mw,bd\})$ & $cd/(spacing+1)$ \\ \cline{2-3}
 & $HPWL_{crit}/(cs+1)$      & $\log(1+HPWL_{ff})$ \\ \cline{2-3}
 & $Asymm_{crit}\!\times\!mw$ & $\rho_{crit}\!\times\!cs$ \\ \cline{2-3}
 & $die\_w/die\_h$            & $cx_{offset}\!\times\!cd$ \\ \cline{2-3}
 & $Eccentricity_{crit}$      & $\log(1+dist_{max}/(cd+1))$ \\ \cline{2-3}
 & $\rho_{crit}/\rho_{global}$ & $\mu_y/die\_h$ \\ \hline
\multirow{16}{*}{\textbf{\shortstack{Wirelength\\(75)}}}
 & \multicolumn{2}{c|}{\textit{Layout Geometry (10)}} \\ \cline{2-3}
 & $\log(1+n_{ff})$; $\log(1+\text{Area})$ & $die\_w/die\_h$; AR \\ \cline{2-3}
 & \multicolumn{2}{c|}{\textit{Cell Composition (11)}} \\ \cline{2-3}
 & $n_{xor}/n_{act}$; $n_{mux}/n_{act}$ & $n_{ff}/n_{act}$; $n_{buf\_inv}/n_{act}$ \\ \cline{2-3}
 & \multicolumn{2}{c|}{\textit{Switching Activity (8)}} \\ \cline{2-3}
 & $\log(1+n_{act}\!\times\!\mu_{ds})$ & $\alpha_{mean}/\alpha_{max}$ \\ \cline{2-3}
 & \multicolumn{2}{c|}{\textit{Timing \& CTS Knobs (12)}} \\ \cline{2-3}
 & $1/t_{clk}$; $\text{Util}\%$ & $cd, cs, mw, bd$ \\ \cline{2-3}
 & \multicolumn{2}{c|}{\textit{Knob$\times$Circuit Interactions (12)}} \\ \cline{2-3}
 & $\log(1+cd\!\times\!\rho_{ff})$ & $\log(1+mw\!\times\!HPWL)$ \\ \cline{2-3}
 & \multicolumn{2}{c|}{\textit{Gravity / Topology Logic-Pull (12)}} \\ \cline{2-3}
 & $\mu(\|\vec{g}\|)$; $\sigma(\|\vec{g}\|)$ & $CV(\|\vec{g}\|)$; $Gini(\|\vec{g}\|)$ \\ \hline
\end{tabular}%
}
\end{table}

\subsection{Independent Prediction Heads}

SwiftCTS employs three independent gradient-boosted prediction heads rather
than a shared-trunk architecture. Because gradient-boosted trees optimize
variance statistics of a single target at a time, a shared trunk would bias
split selection toward objectives with larger variance, reducing sensitivity
to other metrics.

\paragraph{Power head.}
An XGBoost ensemble~\cite{xgboost} on a 20-dimensional feature space.
The target is the scale-free log-ratio $y_\text{power} =
\log(\text{power}/\text{pw\_norm})$, where
$\text{pw\_norm} = n_{ff}\!\times\!f_\text{GHz}\!\times\!\mu_{ds}$
mirrors the dynamic power equation $P = \alpha C_\text{total} V^2 f$,
ensuring the model learns the 10--30\% within-design variation rather than
memorizing per-design absolute baselines.

\paragraph{Wirelength head.}
A blended LightGBM~\cite{lightgbm}--Ridge Regression~\cite{ridgeReg}
architecture (30/70 split) on a 75-dimensional feature space.
The target is $y_\text{wl} = \log(\text{WL}_\text{actual}/\text{wl\_norm})$,
where $\text{wl\_norm} \propto \sqrt{n_{ff}\cdot A}$ approximates the
Steiner-tree minimum~\cite{steinerTree}. A data-driven boundary safeguard
disables the Ridge component when its predicted penalty exceeds the
training-set bounds, preventing numerical explosion on very large OOD designs.

\paragraph{Skew head.}
Clock skew is driven by outliers and lacks a cross-design analytical
normalizer. We therefore reformulate the objective as a
per-placement z-score:
$y_z = (y_\text{ns} - \mu) / \sigma$,
where $\mu$ is the mean skew across historical runs of that specific
placement. A pure LightGBM model on 15 geometric-knob interaction features
predicts this z-score; absolute ns are recovered post-calibration.

\subsection{K-Shot Multiplicative Calibration}
\label{sec:kshot}

Zero-shot predictions on unseen designs achieve strong relative accuracy but
suffer from a constant scalar offset caused by the absolute size and switching
characteristics of the new design. We correct this with a $K$-shot
multiplicative calibration mechanism.
Given $K$ actual CTS runs on a new placement, we compute a geometric-mean
scaling factor to avoid positive bias from arithmetic averaging:
\begin{equation}
  k_\text{cal} = \exp\!\left(\frac{1}{K}\sum_{i=1}^{K}
    \log\!\left(\frac{y_\text{true}^{(i)}}{\hat{y}_\text{pred}^{(i)}}\right)\right)
\end{equation}
This scalar is broadcast across all subsequent predictions for that placement:
$\hat{y}_\text{calibrated} = k_\text{cal}\times\hat{y}_\text{pred}$.
Because the log-ratio formulation has already absorbed cross-design scale
variation, $k_\text{cal}$ corrects only the residual constant offset.
A single reference run ($K=1$) suffices in practice.

\subsection{Search-Based Pareto Optimization}

For design space exploration, every feature derived from the layout geometry
and switching activity is constant for a given placement. We isolate the four
CTS knobs as the only dynamic variables, construct the batch feature matrix
by tiling a single base layout vector $N$ times and patching knob-dependent
columns via vectorized operations, and feed the result directly into the
ensemble inference engines.
The predicted cost matrix $\mathbf{C}\in\mathbb{R}^{N\times 3}$ is normalized
to $[0,1]$, and candidate $i$ is dominated by $j$ if:
$\forall k:\; C_{j,k}\leq C_{i,k}+\varepsilon$ and
$\exists k:\; C_{j,k}<C_{i,k}-\varepsilon$.
We integrate NSGA-II as the search engine (100K evaluations per run,
average 8.9\,s wall clock), but the architecture is algorithm-agnostic.

%%────────────────────────────────────────────────────────────────────────────
\section{GNN Benchmarking Experiments}
\label{sec:gnn}
%%────────────────────────────────────────────────────────────────────────────

\subsection{Experimental Setup}

We evaluate three standard GNN backbones---GCN~\cite{kipf2017semisupervisedclassificationgraphconvolutional},
GraphSAGE~\cite{hamilton2018inductiverepresentationlearninglarge}, and
GATv2~\cite{brody2022attentivegraphattentionnetworks}---on both Raw and
Clustered graph representations.
Each backbone is paired with a Multi-Modal Fusion Architecture that integrates
graph embeddings (via global mean pooling) with scalar placement and CTS knobs
through separate MLPs before a shared regression head predicting three
post-CTS metrics.

\paragraph{Dataset splits.}
\textit{In-distribution (seen):} Four architectures (PicoRV32, AES, SHA-256,
EthMAC), split 80/20 into training and validation.
\textit{Out-of-distribution (unseen):} ZipDiv held out entirely (500 runs)
as a zero-shot generalization test.

All experiments run on an NVIDIA GeForce RTX 5060 (8\,GB VRAM, CUDA~12.8)
for 100 epochs per configuration with batch size 32.

\subsection{Efficiency Results}

Training GATv2 on Raw graphs requires 1.2--2.5\,GB VRAM per batch, already
taxing moderate-scale partitions. The Clustered representation reduces mean
node count by $\sim$$13.3\times$, yielding an average \textbf{17.2$\times$
reduction in peak VRAM} and a \textbf{3$\times$ acceleration in training
throughput}---confirming clustering as a viable strategy on
resource-constrained hardware.

\subsection{Accuracy Trade-off: Global vs.\ Local Metrics}

\paragraph{Global metrics.}
Total Power and Wirelength remain predictable under compression. The Clustered
GCN achieves Power MAE of 0.10 and Wirelength MAE of 0.10 on seen designs,
comparable to the Raw baseline (0.06 each), suggesting clustered proxies are
viable for global objective estimation.

\paragraph{Local metrics.}
Clock Skew prediction degrades sharply under compression.
While Raw GCN achieves Skew MAE of 0.16 with $R^2\approx0.90$ on seen data,
the Clustered models collapse ($R^2 < 0$) on both seen and unseen designs.
On ZipDiv (OOD), the Clustered GCN reaches $R^2 = -2.23$---meaning
predictions are statistically worse than the mean baseline.

\begin{table}[h]
\centering
\small
\caption{GNN Results: MAE and $R^2$ on Raw vs.\ Clustered Representations.}
\label{tab:gnn_results}
\begin{tabular}{l cc cc cc c}
\toprule
\multirow{2}{*}{\textbf{Net}} &
  \multicolumn{2}{c}{\textbf{Skew MAE}} &
  \multicolumn{2}{c}{\textbf{Power MAE}} &
  \multicolumn{2}{c}{\textbf{WL MAE}} &
  \textbf{$R^2$ (S/U)} \\
 & S & U & S & U & S & U & \\
\midrule
\multicolumn{8}{l}{\textit{Raw Graph}} \\
GCN   & 0.16 & 0.24 & 0.06 & 0.07 & 0.06 & 0.11 & 0.95/$-$0.17 \\
SAGE  & 0.16 & 0.26 & 0.06 & 0.06 & 0.06 & 0.11 & 0.95/0.05 \\
GATv2 & 0.16 & 0.22 & 0.06 & 0.06 & 0.06 & 0.11 & 0.95/0.08 \\
\midrule
\multicolumn{8}{l}{\textit{Clustered Graph}} \\
GCN   & 0.17 & 0.16 & 0.10 & 0.13 & 0.10 & 0.15 & 0.85/$-$2.23 \\
SAGE  & 0.17 & 0.12 & 0.09 & 0.09 & 0.09 & 0.16 & 0.90/$-$0.91 \\
GATv2 & 0.17 & 0.06 & 0.09 & 0.13 & 0.09 & 0.23 & 0.89/$-$2.84 \\
\bottomrule
\multicolumn{8}{p{0.44\textwidth}}{\small S = Seen (in-distribution);
  U = Unseen (ZipDiv OOD). $R^2$ is spatial fidelity.}
\end{tabular}
\end{table}

\paragraph{Key finding.}
Generic graph clustering can fundamentally compromise CTS learning objectives
even when global physical metrics remain accurately predicted. This confirms
that future graph coarsening strategies for EDA must be \emph{task-aware},
explicitly preserving structural properties relevant to the downstream
objective (e.g., flop-to-flop path geometry for skew).

%%────────────────────────────────────────────────────────────────────────────
\section{SwiftCTS: Experimental Results}
\label{sec:surrogate_results}
%%────────────────────────────────────────────────────────────────────────────

\subsection{Dataset and Evaluation Protocol}

\paragraph{Dataset.}
All SwiftCTS experiments use CTS-Bench data: 5,520 evaluations from six
designs (four training: AES, PicoRV32, SHA-256, EthMAC; two OOD test:
JPEG encoder, ZipDiv), with 10 CTS configurations per placement.
Training on the larger 19-design corpus (enabled by HPC) will be released
as a follow-on experiment in the camera-ready version.

\paragraph{Evaluation.}
We report MAPE for power and wirelength (normalizes for scale differences
across designs) and MAE in nanoseconds for skew (strict absolute bound on
worst-case temporal error). All predictions are validated against actual
OpenROAD runs.

\paragraph{Two-stage protocol.}
(1)~\textit{Leave-One-Design-Out (LODO)}: trained on three designs,
zero-shot on the fourth --- proves generalization from three designs without
overfitting to specific logic paths.
(2)~\textit{Out-of-Distribution (OOD)}: trained on all four base designs,
deployed zero-shot on JPEG (5$\times$ larger) and ZipDiv (21$\times$ smaller)
--- tests absolute structural limits.

\subsection{Prediction Accuracy}

\begin{table*}[t]
\centering
\small
\renewcommand{\arraystretch}{1.15}
\setlength{\tabcolsep}{5pt}
\caption{SwiftCTS Prediction Results. MAPE (\%) for Power and Wirelength; MAE (ns) for Skew.}
\label{tab:results}
\begin{tabular}{|c|c|rrrr|rrrr|rrrr|}
\hline
\multirow{2}{*}{\textbf{Mode}} &
\multirow{2}{*}{\textbf{Design}} &
\multicolumn{4}{c|}{\textbf{Power (MAPE\%)}} &
\multicolumn{4}{c|}{\textbf{Wirelength (MAPE\%)}} &
\multicolumn{4}{c|}{\textbf{Skew (MAE ns)}} \\
& & K=0 & K=1 & K=2 & K=5 & K=0 & K=1 & K=2 & K=5 &
    K=0 & K=1 & K=2 & K=5 \\ \hline
\multirow{5}{*}{\textbf{LODO}}
 & AES       & 21.0 & 2.2 & 2.0 & \textbf{1.8} & 17.6 & 0.7 & 0.6 & \textbf{0.5}  & -- & 0.1525 & 0.1174 & \textbf{0.0927} \\
 & EthMAC    &  5.7 & 5.0 & 4.3 & \textbf{4.0} &  8.0 & 0.8 & 0.7 & \textbf{0.7}  & -- & 0.1198 & 0.0986 & \textbf{0.0834} \\
 & \cellcolor{bestrow}PicoRV32 & \cellcolor{bestrow}16.8 & \cellcolor{bestrow}2.3 & \cellcolor{bestrow}2.1 & \cellcolor{bestrow}\textbf{1.8}
   & \cellcolor{bestrow}8.0 & \cellcolor{bestrow}0.4 & \cellcolor{bestrow}0.4 & \cellcolor{bestrow}\textbf{0.3}
   & \cellcolor{bestrow}-- & \cellcolor{bestrow}0.0953 & \cellcolor{bestrow}0.0802 & \cellcolor{bestrow}\textbf{0.0665} \\
 & SHA-256   & 54.4 & 3.9 & 3.4 & \textbf{3.2} &  7.5 & 0.4 & \textbf{0.3} & \textbf{0.3} & -- & 0.0982 & 0.0853 & \textbf{0.0713} \\ \cline{2-14}
 & \textbf{Mean} & 24.5 & 3.3 & 2.9 & \textbf{2.7} & 10.3 & 0.6 & 0.5 & \textbf{0.5} & -- & 0.1164 & 0.0954 & \textbf{0.0785} \\ \hline
\multirow{7}{*}{\textbf{OOD (L)}}
 & JPEG-1    & 23.9 & \textbf{3.6} & 4.1 & 4.5 & 22.0 & 2.0 & 1.0 & \textbf{0.7} & -- & 0.2072 & 0.2604 & \textbf{0.2020} \\
 & JPEG-2    & 44.1 & \textbf{6.5} & 7.9 & 7.3 &  7.8 & \textbf{1.1} & 1.2 & \textbf{1.1} & -- & \textbf{0.1366} & 0.1595 & 0.1385 \\
 & \cellcolor{bestrow}JPEG-3 & \cellcolor{bestrow}19.9 & \cellcolor{bestrow}3.6 & \cellcolor{bestrow}\textbf{3.4} & \cellcolor{bestrow}4.7
   & \cellcolor{bestrow}0.8 & \cellcolor{bestrow}1.0 & \cellcolor{bestrow}\textbf{0.8} & \cellcolor{bestrow}\textbf{0.8}
   & \cellcolor{bestrow}-- & \cellcolor{bestrow}0.1149 & \cellcolor{bestrow}0.1208 & \cellcolor{bestrow}\textbf{0.1020} \\
 & JPEG-4    & 37.6 & 3.3 & \textbf{1.6} & 1.7 &  8.3 & \textbf{1.1} & 1.2 & \textbf{1.1} & -- & 0.5434 & 0.2692 & \textbf{0.1849} \\
 & JPEG-5    & 41.8 & 5.3 & \textbf{2.8} & 4.5 &  1.0 & 1.8 & \textbf{0.5} & 0.6 & -- & 0.2648 & 0.0946 & \textbf{0.0937} \\
 & JPEG-6    & 19.9 & \textbf{4.0} & 4.7 & \textbf{4.0} & 16.8 & 1.6 & 2.0 & \textbf{1.1} & -- & 0.2258 & \textbf{0.1652} & 0.1995 \\ \cline{2-14}
 & \textbf{Mean} & 31.2 & 4.4 & \textbf{4.1} & 4.5 &  9.4 & 1.4 & 1.1 & \textbf{0.9} & -- & 0.2488 & 0.1783 & \textbf{0.1534} \\ \hline
\multirow{7}{*}{\textbf{OOD (S)}}
 & ZipDiv-1  &  7.2 & \textbf{5.5} & 6.5 & 7.9 & 63.6 & 0.8 & 0.9 & \textbf{1.1} & -- & \textbf{0.0015} & 0.0016 & 0.0019 \\
 & ZipDiv-2  & 10.9 & 6.4 & \textbf{4.5} & 5.5 & 35.6 & 0.7 & 0.7 & \textbf{0.6} & -- & 0.0038 & 0.0030 & \textbf{0.0018} \\
 & ZipDiv-3  &  9.7 & 2.3 & 3.2 & \textbf{0.5} & 46.3 & 0.6 & 0.3 & \textbf{0.1} & -- & 0.0048 & \textbf{0.0035} & 0.0042 \\
 & \cellcolor{bestrow}ZipDiv-4 & \cellcolor{bestrow}4.6 & \cellcolor{bestrow}\textbf{2.3} & \cellcolor{bestrow}2.8 & \cellcolor{bestrow}2.8
   & \cellcolor{bestrow}60.4 & \cellcolor{bestrow}\textbf{0.4} & \cellcolor{bestrow}0.6 & \cellcolor{bestrow}\textbf{0.4}
   & \cellcolor{bestrow}-- & \cellcolor{bestrow}0.0024 & \cellcolor{bestrow}0.0023 & \cellcolor{bestrow}\textbf{0.0018} \\
 & ZipDiv-5  & 10.6 & 6.8 & 9.7 & \textbf{4.3} & 70.0 & 1.2 & 1.3 & \textbf{0.7} & -- & 0.0066 & 0.0037 & \textbf{0.0026} \\
 & ZipDiv-6  &  7.4 & \textbf{4.7} & 8.0 & 6.5 & 63.6 & \textbf{0.5} & 0.7 & 0.6 & -- & \textbf{0.0040} & 0.0043 & 0.0055 \\ \cline{2-14}
 & \textbf{Mean} & 8.4 & 4.7 & 5.8 & \textbf{4.6} & 56.6 & 0.7 & 0.7 & \textbf{0.6} & -- & 0.0038 & 0.0031 & \textbf{0.0030} \\ \hline
\end{tabular}
\vspace{1mm}
\begin{minipage}{\linewidth}
\small OOD (L) = large OOD designs (JPEG encoder, 5$\times$ larger than
  largest training design); OOD (S) = small OOD designs (ZipDiv, 21$\times$
  smaller). K=0 skew yields relative z-scores; $K\geq1$ required for
  absolute ns.
\end{minipage}
\end{table*}

\noindent
Under zero-shot conditions, SwiftCTS captures relative knob sensitivities
but suffers from scale offsets---ZipDiv exhibits a 56.6\% zero-shot
wirelength MAPE due to severe domain shift.
A single calibration run ($K=1$) reduces LODO power error from 24.5\% to
3.3\% and wirelength error to 0.6\%.
On JPEG (OOD, large), $K=1$ drops wirelength MAPE to 1.4\%.
Timing skew predictions on ZipDiv reach within 5\,ps of physical routing.

\subsection{Design Space Exploration Runtime}

\begin{table}[h]
\centering
\small
\caption{DSE Runtime at 100,000 Evaluations (SwiftCTS).}
\label{tab:runtime}
\begin{tabular}{lll}
\toprule
\textbf{Search Algorithm} & \textbf{Profile} & \textbf{Runtime (avg)} \\
\midrule
Random Search  & Exhaustive  & 604.3\,s \\
Sobol Sequence & Exhaustive  & 600.0\,s \\
\textbf{NSGA-II} & \textbf{Evolutionary} & \textbf{8.9\,s} \\
\bottomrule
\end{tabular}
\end{table}

\noindent
The sub-10-second search budget allows engineers to launch dozens of
differently-seeded NSGA-II runs in parallel within the time of a single
conventional random search iteration, ensuring robust global Pareto coverage.

\subsection{Closed-Loop Physical Validation}

To verify that surrogate-identified configurations improve actual physical
outcomes, we select top Pareto candidates per design and execute them through
the full OpenROAD flow.
On the JPEG encoder (completely unseen OOD design), the actual routed power
(188.80\,mW) matches the surrogate prediction (189.63\,mW) with 0.43\% error.
Clock skew predictions for ZipDiv are within 5\,ps of actual STA.
Across all benchmarks, SwiftCTS-recommended configurations outperform the
OpenROAD default CTS heuristics on all three objectives.

\subsection{Comparison to Prior Frameworks}

\begin{table}[h]
\centering
\small
\caption{ML-Driven CTS Framework Comparison.}
\label{tab:comparison}
\resizebox{\columnwidth}{!}{%
\begin{tabular}{lccccc}
\toprule
\textbf{Framework} & \textbf{GPU?} & \textbf{Train Time} & \textbf{DSE Scale}
  & \textbf{Power} & \textbf{WL} \\
\midrule
Kwon et al.~\cite{kwon2018} (ANN) & Yes & N/R & None & N/R & 13--17\%$^\ast$ \\
Nagaria et al.~\cite{nagaria2020} (CNN) & Yes & N/R & 1-Shot LP & 14.89\% & 8--15\%$^\ast$ \\
GAN-CTS~\cite{gan-cts} & Yes & $\sim$6\,hr & Prob.\ 1-Shot & 1.26\% & 1.64\% \\
\textbf{SwiftCTS (GBDT)} & \textbf{No} & \textbf{4.1\,s} & \textbf{100K in 8.9\,s}
  & \textbf{3.3\%} & \textbf{0.6\%} \\
\bottomrule
\multicolumn{6}{p{0.44\textwidth}}{\small N/R: Not reported.
$^\ast$Intermediate component estimation, not final routed metric.}
\end{tabular}%
}
\end{table}

%%────────────────────────────────────────────────────────────────────────────
\section{Limitations and Future Directions}
\label{sec:limitations}
%%────────────────────────────────────────────────────────────────────────────

\paragraph{Technology node.}
CTS-Bench currently targets Sky130 (130\,nm). Real industrial designs target
5--7\,nm processes with fundamentally different physical constraints.
We plan to extend CTS-Bench to the ASAP7 PDK and eventually to commercial
nodes via anonymized structures.

\paragraph{Surrogate feature space.}
The current 110-dimensional feature space does not model \textit{dont-touch}
nets, non-uniform power grids, or multi-threshold-voltage (multi-\textit{V}t)
cell libraries, all of which influence real clock networks.

\paragraph{Scalability of pipeline.}
The SAIF generation step (gate-level simulation) dominates runtime for
large designs. Future work will investigate faster switching activity
estimation via statistical or RTL-level proxies.

\paragraph{GNN task-aware clustering.}
The GNN benchmarking results (Section~\ref{sec:gnn}) show that generic
graph coarsening breaks skew prediction. Designing clustering objectives
that explicitly preserve flip-flop-to-flip-flop path geometry is a critical
open problem CTS-Bench is well-positioned to drive.

\paragraph{Cross-tool generalization.}
SwiftCTS was validated on OpenROAD. Robustness under commercial tools
(Cadence Innovus, Synopsys ICC2) requires industrial collaboration and
remains future work.

%%────────────────────────────────────────────────────────────────────────────
\section{Conclusion}
\label{sec:conclusion}
%%────────────────────────────────────────────────────────────────────────────

We presented CTS-Bench, an open, large-scale dataset and automated pipeline
for AI-driven Clock Tree Synthesis.
By combining 10,000+ CTS evaluations across 19 diverse RTL designs, a
production-grade HPC generation pipeline, multi-scale graph representations,
and the SwiftCTS surrogate as a concrete demonstration of dataset utility,
CTS-Bench addresses the core data bottleneck that has constrained ML research
on CTS for over a decade.

The dataset enables two classes of research that were previously impractical:
(1)~GNN benchmarking that honestly characterizes the accuracy-efficiency
trade-off of graph coarsening for physical design; and (2)~surrogate-based
design space exploration that generalizes across unseen macro architectures
with minimal physical reference data.
All data, code, and the HPC pipeline are publicly released to encourage
community adoption and extension.

\bibliographystyle{ACM-Reference-Format}
\bibliography{todaes-refs}

\end{document}
